\documentclass[12pt,a4paper,onecolumn]{exam}

\usepackage{amsmath}
\usepackage{amsthm}
\usepackage{amssymb}
\usepackage{graphicx}
\usepackage{float}
\usepackage{geometry}
\usepackage{tikz}
\usepackage[skins]{tcolorbox}
\usepackage{listings}
\usepackage{xcolor}
\usepackage{hyperref}
\usepackage{multirow}
\usepackage{adjustbox}
\usepackage{booktabs} 
\usepackage[font=small]{caption}
\usepackage{subcaption}

% --- Code Listing Colors ---
\definecolor{codegreen}{rgb}{0,0.6,0}
\definecolor{codegray}{rgb}{0.5,0.5,0.5}
\definecolor{codepurple}{rgb}{0.58,0,0.82}
\definecolor{backcolour}{rgb}{0.95,0.95,0.92}

% --- Code Listing Style ---
\lstdefinestyle{mystyle}{
    backgroundcolor=\color{backcolour},
    commentstyle=\color{codegreen},
    keywordstyle=\color{magenta},
    numberstyle=\tiny\color{codegray},
    stringstyle=\color{codepurple},
    basicstyle=\ttfamily\footnotesize,
    breaklines=true,
    captionpos=b,
    keepspaces=true,
    numbers=left,
    numbersep=2pt,
    showspaces=false,
    showstringspaces=false,
    showtabs=false,
    tabsize=1
}

\lstset{style=mystyle, language=Python}

\begin{document}

\renewcommand{\qedsymbol}{$\blacksquare$}
\begingroup
    \centering
    \LARGE \textbf{E1 251 Course Project}\\
    \LARGE \textbf{Reconstruction from Non-Uniform Samples Using a DCT-$l_p$ Prior(MM-ADMM)}\\
    \phantom{123} \\
    \large Dwaipayan Haldar(SR. No.: 27128)\par
\endgroup

\noindent\rule{\textwidth}{0.5pt}

\section{Description of Experimental setup}
\label{section:Experimental Setup}

\subsection{Model Images}
Images used for the project are \verb|cameraman.tif| and \verb|lena|. Both have a resolution of $256 \times 256$ and normalized to $[0,1]$ using \verb|img_as_float| function of \verb|skimage| package.

\subsection{Sampling and Noise}
I have used binary random masks using the \verb|numpy.random| with $r \in \{0.2,0.3,0.5\}$, where $r$ is the fraction of pixels used. \\
Noise $\eta_i \sim \mathcal{N}(0,\sigma^2)$, where $\sigma$ corresponds to an SNR of 30dB. A closed form value of $\sigma$ can be derived as follows:
\[
\begin{aligned}
\text{SNR} = 20\log_{10}\frac{||Wx^\star||_2}{||\eta||_2} = 30
\implies \frac{||Wx^\star||_2}{||\eta||_2} = 10^{1.5} \implies ||\eta||_2^2 = \frac{||Wx^\star||_2^2}{10^{3}}
\end{aligned}
\]
$\because \mathbb{E}(\eta_i^2) = \sigma^2$(Normal Distribution with 0 Mean)
\[
\implies ||\eta||_2^2 = N\sigma^2 \implies \sigma^2 = \frac{||Wx^\star||_2^2}{M.10^{3}} \implies \sigma = \sqrt{\frac{||Wx^\star||_2^2}{M.10^{3}}}
\]
This is the formula used to calculate the random noise in \ref{code:samplingfunction}. 

\subsection{Choice of parameters, Evaluation and Experiments}
Experiments are run for each choice of $(r,p)$. Results are shown in \ref{section:results} for best $\lambda$ in case of each $(r,p)$. Best $\lambda$ is choosen on the basis of the highest PSNR value. Table.\ref{table:resultstable} gives the concise report for both the images.


\section{Results}
\label{section:results}

\subsection{Cameraman}
Fig.\ref{img:cameraman0.2} gives the original image, sampled noisy image and the Reconstructed image for sampling percentage 0.2. Similarily, Fig.\ref{img:cameraman0.3} and Fig.\ref{img:cameraman0.5} gives the same results for sampling percentage 0.3 and 0.5 respectively. Fig.\ref{img:cameraman0.5psnrlambda} gives the PSNR vs $\lambda$ plot for r = 0.5 and p $\in \{0.5,0.7,0.9\}$. Fig.\ref{img:cameraman0.5p0.53plots} gives the Objective Function vs iteration step, No. of ADMM iteration vs iteration step, Relative error vs iteration step plots for r = 0.5, p = 0.3 and $\lambda = 0.01$ which is the best $\lambda$ for the given configuration. Qualitatively, p = 0.5 outperforms p = 0.4, 0.3 which would evident from Table.\ref{table:resultstable}, also for r = 0.5 we get better results than r=0.2, 0.3 which is quite expected. 
\begin{figure}[H]
    \centering
    \includegraphics[scale = 0.3]{output images/Cameraman_0.2.png}
    \caption{(First Row)Original Image, Sampling Mask, Noisy Masked Image (Second Row)Reconstructed Image(p = 0.5,0.7,0.9) for \textbf{r = 0.2}}
    \label{img:cameraman0.2}
\end{figure}
\begin{figure}[H]
    \centering
    \includegraphics[scale = 0.3]{output images/Cameraman_0.3.png}
    \caption{(First Row)Original Image, Sampling Mask, Noisy Masked Image (Second Row)Reconstructed Image(p = 0.5,0.7,0.9) for \textbf{r = 0.3}}
    \label{img:cameraman0.3}
\end{figure}
\begin{figure}[H]
    \centering
    \includegraphics[scale = 0.3]{output images/Cameraman_0.5.png}
    \caption{(First Row)Original Image, Sampling Mask, Noisy Masked Image (Second Row)Reconstructed Image(p = 0.5,0.7,0.9) for \textbf{r = 0.5}}
    \label{img:cameraman0.5}
\end{figure}
\begin{figure}[H]
    \centering
    \includegraphics[scale = 0.35]{output images/Cameraman_0.5_psnr_lambda.png}
    \caption{PSNR vs $\lambda$ for 3 combination of (r,p) where r = 0.5 and p $\in \{0.5,0.7,0.9\}$}
    \label{img:cameraman0.5psnrlambda}
\end{figure}
\begin{figure}[H]
    \centering
    \includegraphics[scale = 0.35]{output images/Cameraman_0.5_p_0.3_3plots.png}
    \caption{(From left) Objective Function vs iteration step, No. of ADMM iteration vs iteration step, Relative error vs iteration step for (r,p,$\lambda$(r,p))=(0.5, 0.3, 0.01)} 
    \label{img:cameraman0.5p0.53plots}
\end{figure}

\subsection{Lena}
Fig.\ref{img:lena0.2} gives the original image, sampled noisy image and the Reconstructed image for sampling percentage 0.2. Similarily, Fig.\ref{img:lena0.3} and Fig.\ref{img:lena0.5} gives the same results for sampling percentage 0.3 and 0.5 respectively. Fig.\ref{img:lena0.5psnrlambda} gives the PSNR vs $\lambda$ plot for r = 0.5 and p $\in \{0.5,0.7,0.9\}$. Fig.\ref{img:lena0.5p0.53plots} gives the Objective Function vs iteration step, No. of ADMM iteration vs iteration step, Relative error vs iteration step plots for r = 0.5, p = 0.3 and $\lambda = 0.01$ which is the best $\lambda$ for the given configuration. Similar results are seen for lena image also. Qualitatively, p = 0.5 outperforms p = 0.4, 0.3 which is also supported from Table.\ref{table:resultstable}, also for r = 0.5 we get better results than r=0.2, 0.3 which is quite expected. 
\begin{figure}[H]
    \centering
    \includegraphics[scale = 0.3]{output images/lena_0.2.png}
    \caption{(First Row)Original Image, Sampling Mask, Noisy Masked Image (Second Row)Reconstructed Image(p = 0.5,0.7,0.9) for \textbf{r = 0.2}}
    \label{img:lena0.2}
\end{figure}
\begin{figure}[H]
    \centering
    \includegraphics[scale = 0.3]{output images/lena_0.3.png}
    \caption{(First Row)Original Image, Sampling Mask, Noisy Masked Image (Second Row)Reconstructed Image(p = 0.5,0.7,0.9) for \textbf{r = 0.3}}
    \label{img:lena0.3}
\end{figure}
\begin{figure}[H]
    \centering
    \includegraphics[scale = 0.3]{output images/lena_0.5.png}
    \caption{(First Row)Original Image, Sampling Mask, Noisy Masked Image (Second Row)Reconstructed Image(p = 0.5,0.7,0.9) for \textbf{r = 0.5}}
    \label{img:lena0.5}
\end{figure}
\begin{figure}[H]
    \centering
    \includegraphics[scale = 0.35]{output images/lena_0.5_psnr_lambda.png}
    \caption{PSNR vs $\lambda$ for 3 combination of (r,p) where r = 0.5 and p $\in \{0.5,0.7,0.9\}$}
    \label{img:lena0.5psnrlambda}
\end{figure}
\begin{figure}[H]
    \centering
    \includegraphics[scale = 0.35]{output images/lena_0.5_p_0.3_3plots.png}
    \caption{(From left) Objective Function vs iteration step, No. of ADMM iteration vs iteration step, Relative error vs iteration step for (r,p,$\lambda$(r,p))=(0.5, 0.3, 0.01)} 
    \label{img:lena0.5p0.53plots}
\end{figure}


\subsection{Combined Results}

\begin{table}[H]
\centering
\caption{Comparison of (r,p) values for both the images}
\label{table:resultstable}
% \tiny
\renewcommand{\arraystretch}{1.1}

\begin{adjustbox}{width=0.9\textwidth}
\begin{tabular}{|c|c|c|c|c|c|c|}
\hline
Image & r & p & $\lambda(r,p)$ & PSNR & Relative Error & Runtime(s)\\
\hline

\multirow{9}{*}{cameraman}


& \multirow{3}{*}{0.2} & 0.5 & 0.1 & 19.81 & 0.194 & 1.92\\ 
&                      & 0.7 & 0.1 & 19.82 & 0.194 & 1.73\\ 
&                      & \textbf{0.9} & 0.01 & \textbf{20.25} & \textbf{0.184} & 11.71 \\ \cline{2-7}


& \multirow{3}{*}{0.3} & 0.5 & 0.01 & 20.78 & 0.173 & 6.57\\ 
&                      & 0.7 & 0.01 & 21.57 & 0.158 & 6.54\\ 
&                      & \textbf{0.9} & 0.01 & \textbf{22.15} & \textbf{0.148} & 10.90\\ \cline{2-7}


& \multirow{3}{*}{0.5} & 0.5 & 0.01 & 23.85 & 0.122 & 7.32\\ 
&                      & 0.7 & 0.01 & 24.51 & 0.113 & 3.8\\ 
&                      & \textbf{0.9} & 0.01 & \textbf{24.93} & \textbf{0.107} & 4.31\\

\hline


\multirow{9}{*}{lena}

& \multirow{3}{*}{0.2} & 0.5 & 0.01 & 21.69 & 0.159 & 7.12\\ 
&                      & 0.7 & 0.01 & 22.52 & 0.144 & 6.72\\ 
&                      & \textbf{0.9} & 0.01 & \textbf{23.09} & \textbf{0.135} & 6.25\\ \cline{2-7}


& \multirow{3}{*}{0.3} & 0.5 & 0.01 & 24.11 & 0.120 & 4.55\\ 
&                      & 0.7 & 0.01 & 24.62 & 0.113 & 7.25\\ 
&                      & \textbf{0.9} & 0.01 & \textbf{24.96} & \textbf{0.109} & 7.64\\ \cline{2-7}


& \multirow{3}{*}{0.5} & 0.5 & 0.01 & 26.97 & 0.086 & 2.67\\ 
&                      & 0.7 & 0.01 & 27.52 & 0.081 & 2.61\\ 
&                      & \textbf{0.9} & 0.01 & \textbf{27.80} & \textbf{0.078} & 3.04 \\

\hline


\end{tabular}
\end{adjustbox}
% \caption{Results for Img1 with varying $(r,p)$ values}
% \label{table:results table}
\end{table}

\section{Discussion}
\label{section:discussion}

\subsection{Effect of \texorpdfstring{$\lambda$}{lambda}}
$\lambda$ is the weight given to the sparsity term. Here $\lambda$ is varied over the log scale $\{10^{-4},10^{-3},10^{-2},10^{-1},1\}$. Almost in all the cases, $\lambda$ = 0.1 or 0.01 is the winner. That is expected since $\lambda = 10^{-4}$ means too less weightage to the sparsity term. So, the images are likely to deviate from natural images. Because, natural images is sparse in the DCT domain. $\lambda = 1$ is giving too much weightage to the sparsity term. So, the image is likely to be oversmoothed. So, a moderate value of $\lambda$ is expected to give a balanced approach. 

\subsection{Effect of p}
p is another parameter that controls the sparsity. As per Table.\ref{table:resultstable}, it is quite evident that p = 0.9 comes to be the winner in all the cases. This is because p = 0.5 is an over aggressive norm for the sparsity term. Whereas p = 0.9 gives the most effective balance for the sparsity prior. 

\subsection{Effect of r}
Understanding the effect of r is pretty trivial. As per Table.\ref{table:resultstable}, with increase in r there is a consistent increase in PSNR, consistent decrease in relative error. r is the fraction of samples selected. It is evident that if more samples are selected the final results will come out to better than less no. of samples selected. 
\phantom{123}\\
\phantom{123}\\
The next section provides the commented code for all the functions used. All the output images are attached in the zip file. It contains all the plots for all possible combinations. \verb|src.py| contains the source code for all the functions. \verb|wrapper.ipynb| contains the outputs of those functions.   

\newpage
\appendix
\section{Code}

\subsection{Main Functions}
This section contains all the main functions used for the project namely the random sampling function, the MM inner loop, the MM-ADMM  complete loop.
\subsubsection{Sampling Function}
\begin{lstlisting}[label=code:samplingfunction]
def sampling(img, r):
    """
    Argument: img and sampling ratio r
    Output: Sampled Noisy Vector m and indices idx
    """
    img_vector = img.flatten() 
    N = img_vector.shape[0] 
    M = int(np.round(r * N)) 
    
    idx = np.random.choice(N, size=M, replace=False) 
    Wx_ = img_vector[idx] 

    # Noise addition
    Wx_l2_2 = (np.linalg.norm(Wx_))**2 
    sigma = np.sqrt(Wx_l2_2 / (1000 * M)) # Approx 30dB SNR
    noise = np.random.normal(0, sigma, (M,)) 
    return (Wx_ + noise), idx 
\end{lstlisting}

\subsubsection{ADMM  Inner loop}
\begin{lstlisting}
def admm_inner_solver(N, m, idx, w_k, x_init, y_prev, lammbda, rho, img_shape, max_admm_iter=1000, tol=1e-4):
    """
    Solves the Weighted L1 subproblem using ADMM.
    """
    x = x_init.copy()
    v = y_prev.copy()
    u = np.zeros_like(v)     
    
    mask = np.zeros(N)
    mask[idx] = 1.0       
    
    m_full = np.zeros(N)
    m_full[idx] = m       

    iteration = 0
    
    for t in range(max_admm_iter):
        v_prev = v.copy()
        
        #  x-update 
        s_t = idct2(v - u, img_shape)
        x = mask * ((m_full + rho * s_t) / (1 + rho)) + (1 - mask) * s_t
        
        #  v-update (DCT Domain)
        dct_x = dct2(x, img_shape)
        z_t = dct_x + u
        
        # Weighted Soft Thresholding
        threshold = (lammbda / rho) * w_k
        v = soft_threshold(z_t, threshold)
        
        #  u-update (Dual) 
        u = u + dct_x - v
        
        #  Convergence Check 
        r_norm = np.linalg.norm(dct_x - v)      # Primal residual
        s_norm = np.linalg.norm(rho * (v - v_prev)) # Dual residual
        
        iteration = t + 1
        if r_norm < tol and s_norm < tol:
            break
            
    return x, v, iteration   
\end{lstlisting}

\subsubsection{MM ADMM  Outer loop}
\begin{lstlisting}
def mm_admm(N, m, idx, rho, lammbda, p, img_shape, max_mm_iter=100, tol=1e-3, epsilon=1e-6):
    # Initialize x with zero-filled observations
    x = np.zeros(N)
    x[idx] = m
    
    # Initial DCT coefficients
    v = dct2(x, img_shape)
    
    list_iteration_admm = [] 
    list_objective_function = [] 
    list_relative_error = [] 
    
    for k in range(max_mm_iter):
        x_prev = x.copy()
        y_prev = v
        
        #  MM Weight Update 
        w_k = p * (np.abs(y_prev) + epsilon)**(p - 1)
            
        #  ADMM Inner Solver 
        x, v, inner_iters = admm_inner_solver(
            N, m, idx, w_k, x_prev, y_prev, lammbda, rho, img_shape
        )
        
        #  Convergence Check 
        norm_x = np.linalg.norm(x)
        rel_error = np.linalg.norm(x - x_prev) / max(1, norm_x)

        # Logging
        obj_val = objective_function(idx, x, m, lammbda, p, img_shape)
        list_objective_function.append(obj_val)
        list_iteration_admm.append(inner_iters)
        list_relative_error.append(rel_error)
        
        if rel_error < tol:
            break
            
    return x, list_objective_function, list_iteration_admm, list_relative_error
\end{lstlisting}

\subsection{Auxiliary Functions}

\subsubsection{Plotting Function}
This is the outermost function, when we call with image and r value, gives all the plots and all the values like PSNR, Relative Error etc. It is the function that is called for output in jupyter notebook.

\begin{lstlisting}
def reconstruct(img, r):
    img_shape = img.shape
    N = img_shape[0] * img_shape[1] 
    
    # Generate Samples
    m, idx = sampling(img, r) 
    
    # For Visualization
    sampling_mask_image = sampling_mask(N, idx).reshape(img_shape)
    noisy_image = Wtranspose(N, m, idx).reshape(img_shape)

    # Storage for results across different p values
    results_by_p = [] # Will store dictionary of results for each p
    
    p_values = [0.5, 0.7, 0.9] # As per Algorithm 1
    
    print(f"Starting Reconstruction for r={r}...")
    start_time_total = time.time()
    
    for p in p_values:
        print(f"\n Processing p={p} ")
        
        # Initialize lists to store sensitivity data for this p
        psnrs_for_plot = []
        lambdas_for_plot = []
        
        best_metrics = {
            "psnr": -np.inf,
            "lambda": None,
            "image": None,
            "obj_hist": [],
            "iter_hist": [],
            "err_hist": []
        }
        
        # Grid search for Lambda
        lambda_grid = np.logspace(-4, 0, num=5) 
        
        for l in lambda_grid:
            start_time_lambda = time.time()
            
            # Run the algorithm
            x_rec_flat, obj_hist, iter_hist, err_hist = mm_admm(
                N, m, idx, rho=4.0, lammbda=l, p=p, img_shape=img_shape
            )
            
            end_time_lambda = time.time()
            runtime_lambda = end_time_lambda - start_time_lambda
            
            # Reshape and compute metrics
            x_rec = x_rec_flat.reshape(img_shape)
            curr_psnr = psnr(x_rec, img)
            curr_rel_error = relative_change(x_rec, img)
            
            # Store data for plotting
            psnrs_for_plot.append(curr_psnr)
            lambdas_for_plot.append(l)
            
            # Update best metrics if this is the best PSNR so far
            if curr_psnr > best_metrics["psnr"]:
                best_metrics["psnr"] = curr_psnr
                best_metrics["lambda"] = l
                best_metrics["image"] = x_rec
                best_metrics["obj_hist"] = obj_hist
                best_metrics["iter_hist"] = iter_hist
                best_metrics["err_hist"] = err_hist
                best_metrics["rel_error"] = curr_rel_error
                best_metrics["runtime"] = runtime_lambda

        # Save the sensitivity lists into the dictionary
        best_metrics["psnr_list_for_plot"] = psnrs_for_plot
        best_metrics["lambda_list_for_plot"] = lambdas_for_plot
        
        results_by_p.append(best_metrics)
        print(f"Best for p={p}: Lambda={best_metrics['lambda']:.1e}, PSNR={best_metrics['psnr']:.2f}")
        print(f"Runtime for best lambda: {best_metrics['runtime']:.4f}s, Relative error: {best_metrics['rel_error']:.4e}")


    end_time_total = time.time()
    print(f"\nTotal runtime: {end_time_total - start_time_total:.2f}s")

    #  Plotting 
    
    # 1. Visual Comparison
    plt.figure(figsize=(15, 10))
    
    plt.subplot(2, 3, 1)
    plt.title("Original Image")
    plt.imshow(img, cmap="gray")
    plt.axis('off')
    
    plt.subplot(2, 3, 2)
    plt.title("Sampling Mask")
    plt.imshow(sampling_mask_image, cmap="gray")
    plt.axis('off')
    
    plt.subplot(2, 3, 3)
    plt.title(f"Noisy Obs (r={r})")
    plt.imshow(noisy_image, cmap="gray")
    plt.axis('off')
    
    for i, p in enumerate(p_values):
        plt.subplot(2, 3, 4 + i)
        plt.title(f"Rec (p={p})\nPSNR: {results_by_p[i]['psnr']:.2f}dB")
        plt.imshow(results_by_p[i]['image'], cmap="gray")
        plt.axis('off')
    plt.tight_layout()
    plt.show()

    # 2. Convergence Metrics
    for i, p in enumerate(p_values):
        res = results_by_p[i]
        k_range = range(1, len(res['obj_hist']) + 1)
        
        plt.figure(figsize=(18, 4))
        plt.suptitle(f"Convergence Metrics for p={p} (Best Lambda={res['lambda']:.1e})")
        
        plt.subplot(1, 3, 1)
        plt.plot(k_range, res['obj_hist'])
        plt.title("Objective Function vs Outer Iter")
        plt.xlabel("Outer Iter (k)")
        plt.ylabel("J(x)")
        plt.grid(True)
        
        plt.subplot(1, 3, 2)
        plt.plot(k_range, res['iter_hist'])
        plt.title("ADMM Iterations per MM Step")
        plt.xlabel("Outer Iter (k)")
        plt.ylabel("Inner Iters")
        plt.grid(True)
        
        plt.subplot(1, 3, 3)
        plt.plot(k_range, res['err_hist'])
        plt.title("Relative Error vs Outer Iter")
        plt.xlabel("Outer Iter (k)")
        plt.yscale('log')
        plt.ylabel("||x_k+1 - x_k|| / ||x_k||")
        plt.grid(True)
        
        plt.show()

    # 3. Lambda Sensitivity Plot (New)
    plt.figure(figsize=(10, 6))
    for i, p in enumerate(p_values):
        res = results_by_p[i]
        plt.semilogx(
            res["lambda_list_for_plot"], 
            res["psnr_list_for_plot"], 
            'o-', 
            linewidth=2, 
            label=f'p={p}'
        )
    
    plt.title("Parameter Sensitivity: PSNR vs $\lambda$")
    plt.xlabel("$\lambda$ (log scale)")
    plt.ylabel("PSNR (dB)")
    plt.legend()
    plt.grid(True, which="both", ls="-")
    plt.tight_layout()
    plt.show()

    return [res['image'] for res in results_by_p]
\end{lstlisting}

\subsubsection{Other Functions}
This section contains those functions like psnr, sampling mask etc. 
\begin{lstlisting}
def dct2(x_flat, shape):
    """
    Performs 2D Orthonormal DCT on a flattened vector.
    Reshapes to 2D -> Applies DCT -> Flattens back.
    """
    x = x_flat.reshape(shape)
    return dct(dct(x.T, norm='ortho').T, norm='ortho').flatten()

def idct2(x_flat, shape):
    """
    Performs 2D Orthonormal IDCT on a flattened vector.
    """
    x = x_flat.reshape(shape)
    return idct(idct(x.T, norm='ortho').T, norm='ortho').flatten()

def psnr(img, img_ref):
    img = img.flatten()
    img_ref = img_ref.flatten()
    num = 1.0 
    if np.max(img_ref) > 1.1: num = 255.0
    den = np.linalg.norm(img_ref - img) / np.sqrt(img_ref.shape[0])
    if den == 0: return 100
    return 20 * np.log10(num / den)

def relative_change(x, y):
    x = x.flatten()
    y = y.flatten()
    if np.linalg.norm(y) == 0: return 0
    return np.linalg.norm(x - y) / np.linalg.norm(y)

def sampling_mask(N, idx):
    mask = np.zeros((N,))
    mask[idx] = 1
    return mask

def Wtranspose(N, x, idx):
    out = np.zeros((N,))
    out[idx] = x
    return out

def objective_function(idx, x, m, lammbda, p, img_shape, epsilon=1e-6):
    diff = x[idx] - m
    term1 = np.sum(diff**2)
    y = dct2(x, img_shape)
    term2 = lammbda * np.sum((epsilon + np.abs(y))**p)
    return term1 + term2

def soft_threshold(z, threshold):
    return np.sign(z) * np.maximum(np.abs(z) - threshold, 0)
\end{lstlisting}
After this, all the output is obtained by simply calling the \verb|reconstruct| function with the image and the sampling factor in the \verb|wrapper.ipynb| file and the output images are saved in the respective folder.


\end{document}